
% VLDB template version of 2020-08-03 enhances the ACM template, version 1.7.0:
% https://www.acm.org/publications/proceedings-template

\documentclass[sigconf, nonacm]{acmart}

\newcommand\vldbdoi{XX.XX/XXX.XX}
\newcommand\vldbpages{XXX-XXX}
\newcommand\vldbvolume{14}
\newcommand\vldbissue{1}
\newcommand\vldbyear{2020}

\usepackage{algorithm}
\usepackage{algpseudocode}
\usepackage{booktabs}
\usepackage{amsthm}

\newtheorem{theorem}{Theorem}[section]
\newtheorem{lemma}[theorem]{Lemma}
\newtheorem{corollary}[theorem]{Corollary}
\newtheorem{proposition}[theorem]{Proposition}
\theoremstyle{definition}
\newtheorem{definition}[theorem]{Definition}
\theoremstyle{remark}
\newtheorem{remark}[theorem]{Remark}

\newcommand\vldbauthors{\authors}
\newcommand\vldbtitle{\shorttitle}
\newcommand\vldbavailabilityurl{URL_TO_YOUR_ARTIFACTS}
\newcommand\vldbpagestyle{plain}

\begin{document}
\title{Redundancy-Free Maximal Clique Enumeration via Search-Tree Reordering}

\begin{comment}
\author{Author Name}
\affiliation{\institution{Institution}}
\email{author@institution.edu}
\end{comment}

\begin{abstract}
We study exact maximal clique enumeration in an undirected graph.  A naive
search over all vertex subsets is infeasible, and Bron--Kerbosch--style search
greatly improves on this by expanding only clique-consistent candidates.
However, even after a maximal clique $C$ has been reported, the search still
contains branches rooted at proper subsets of $C$.  Those branches are known
not to terminate inside $C$, yet standard depth-first search and pivoting do
not use that information to reorganize the remaining search space.

We present \textsc{ReorderSib}, an exact algorithm based on recursive
search-tree reordering.  The central observation is that once a branch state
$M$ is known to lie inside a discovered maximal clique $C$, every undiscovered
maximal clique extending $M$ must escape $C$ through a vertex in $Q \setminus
C$.  For a single covering clique, this yields an explicit reordering rule:
outside vertices must be explored before inside vertices.  For several
covering cliques, the admissible escape choices become exactly the minimal
clique-compatible hitting sets of the induced constraint family.  These sets
seed a replacement family of branches that preserves completeness while
avoiding redundant traversal of subsets of known cliques.

We formalize this idea, prove soundness, completeness, and duplicate-freedom
for the resulting core algorithm, and then describe how the implementation
realizes the same principle using clique indexing, branch merging, and exact
hitting-set solvers.
\end{abstract}

\maketitle

%%% do not modify the following VLDB block %%
%%% VLDB block start %%%
\pagestyle{\vldbpagestyle}
\begingroup\small\noindent\raggedright\textbf{PVLDB Reference Format:}\\
\vldbauthors. \vldbtitle. PVLDB, \vldbvolume(\vldbissue): \vldbpages, \vldbyear.\\
\href{https://doi.org/\vldbdoi}{doi:\vldbdoi}
\endgroup
\begingroup
\renewcommand\thefootnote{}\footnote{\noindent
This work is licensed under the Creative Commons BY-NC-ND 4.0 International
License. Visit \url{https://creativecommons.org/licenses/by-nc-nd/4.0/} to
view a copy of this license. For any use beyond those covered by this license,
obtain permission by emailing \href{mailto:info@vldb.org}{info@vldb.org}.
Copyright is held by the owner/author(s). Publication rights licensed to the
VLDB Endowment.\\
\raggedright Proceedings of the VLDB Endowment,
Vol.~\vldbvolume, No.~\vldbissue\ ISSN 2150-8097.\\
\href{https://doi.org/\vldbdoi}{doi:\vldbdoi}\\
}\addtocounter{footnote}{-1}\endgroup
%%% VLDB block end %%%

\ifdefempty{\vldbavailabilityurl}{}{
\vspace{.3cm}
\begingroup\small\noindent\raggedright\textbf{PVLDB Artifact Availability:}\\
The source code, data, and/or other artifacts have been made available at
\url{\vldbavailabilityurl}.
\endgroup
}

% ============================================================
\section{Introduction}
\label{sec:introduction}
% ============================================================

A \emph{maximal clique} in an undirected graph $G=(V,E)$ is a complete subgraph
that cannot be extended by any additional vertex.  The problem of enumerating all
maximal cliques is a fundamental primitive in graph mining, with applications in
community detection~\cite{palla2005uncovering}, protein-complex
identification~\cite{bader2003automated}, and network security analysis.

A direct way to solve the problem is to examine subsets of vertices and test
whether each subset is a maximal clique.  This view makes the combinatorial
difficulty clear, but it is also hopeless in practice because there are
$2^{|V|}$ subsets.  Exact algorithms therefore organize the search much more
carefully.  The Bron--Kerbosch algorithm~\cite{bron1973algorithm} and its
pivoting variant~\cite{tomita2006worst} provide the standard framework: they
maintain a growing clique $R$, a set $P$ of candidate extensions, and a set
$X$ of already processed vertices, and branch only on vertices that remain
clique-consistent with the current partial solution.  Combined with the
degeneracy ordering of Eppstein, L{\"o}ffler, and Strash~\cite{eppstein2010listing},
this framework achieves the worst-case optimal bound
$O(\delta \cdot n \cdot 3^{\delta/3})$, where $\delta$ is the degeneracy of
$G$.

\smallskip

A structural redundancy nevertheless remains.  Suppose the search reports a
maximal clique $C = \{a,b,c,d\}$.  At that moment the search tree may still
contain branches rooted at $\{a\}$, $\{a,b\}$, and $\{a,b,c\}$, because those
branches were created before $C$ was known.  None of these states can terminate
inside $C$: any new maximal clique reachable from $\{a,b,c\}$ must contain a
vertex outside $C$, or else it would be a subset of $C$ and therefore not
maximal.  Standard pivoting can delay some branches, but it does not
systematically re-route the search once such information becomes available.

This paper is about exploiting that missing information.  Our goal is not to
replace clique-consistent search, but to refine it: after a clique is found, we
want the remaining search to focus first on the vertices that can lead to
genuinely new maximal cliques and to postpone vertices that are already
explained by the discovered clique.

\smallskip

\noindent\textbf{Our contribution.}\quad We present \textsc{ReorderSib}, an
exact maximal-clique enumeration algorithm based on recursive search-tree
reordering.  Its central mechanism is simple: once a discovered maximal clique
$C$ covers a branch state $M$, the search below $M$ must leave $C$, so vertices
from $Q \setminus C$ are promoted to the front of the search and vertices from
$Q \cap C$ are pushed to the tail.  For one covering clique this gives an
explicit reordering rule; for several covering cliques it leads to a minimal
clique-compatible hitting-set problem whose solutions define a replacement
family of branches.

The paper develops this idea in three stages.  First, we formalize the
structural facts that make reordering correct: canonical ownership of a clique
by a unique root branch, the escape property induced by a discovered clique,
and the sibling-set characterization of valid replacement branches.  Second, we
turn these facts into a clean core algorithm and prove soundness
(Theorem~\ref{thm:sound}), completeness (Theorem~\ref{thm:complete}), and
duplicate-freedom (Theorem~\ref{thm:unique}).  Third, we explain how the actual
implementation realizes the same idea using clique indexing, branch merging,
and exact hitting-set solvers.

The paper is organised as follows.  Section~\ref{sec:prelim} fixes notation.
Section~\ref{sec:dfs} turns the redundancy observed above into precise
structural lemmas.  Section~\ref{sec:algorithm} presents the core algorithm and
its correctness proof.  Section~\ref{sec:implementation} then explains how the
\textsc{ReorderSib} implementation realizes the same search principle in code.

% ============================================================
\section{Preliminaries}
\label{sec:prelim}
% ============================================================

Let $G = (V, E)$ be a simple, undirected graph with $n = |V|$ vertices and
$m = |E|$ edges.  For a vertex $v \in V$, write $N(v)$ for its open
neighborhood.  For a set $S \subseteq V$, write $N(S) = (\bigcap_{v \in S}
N(v)) \setminus S$ for the set of vertices adjacent to every member of $S$ but
not in $S$ itself.  A \emph{clique} is a set $K \subseteq V$ in which every
pair of vertices is adjacent.  A clique is \emph{maximal} if no vertex outside
it is adjacent to all of its members.  We write $\mathcal{M}(G)$ for the family
of all maximal cliques; the Moon--Moser bound gives $|\mathcal{M}(G)| \le
3^{n/3}$~\cite{moon1965cliques}.

\paragraph{Vertex ordering.}
Fix a total order $\sigma \colon V \to [n]$ --- in practice, the degeneracy
order~\cite{eppstein2010listing}, computed in $O(n+m)$ by iteratively removing a
minimum-degree vertex.  Define the \emph{forward neighborhood} of $v$ as
\[
  N^+(v) \;=\; \bigl\{\, u \in N(v) : \sigma(u) > \sigma(v) \,\bigr\}.
\]
Under the degeneracy order, $|N^+(v)| \le \delta$ for every $v$, where $\delta$
is the degeneracy of $G$.

\paragraph{Branches.}
The search state at any point in the algorithm is captured by a \emph{branch}.

\begin{definition}[Branch]\label{def:branch}
A \emph{branch} is a triple $(M, h, Q)$ where $M$ is a clique in $G$,
$h \in M$ is the \emph{branch head}, and $Q \subseteq N(M) \setminus M$ is the
\emph{candidate set}.  We call $M$ the \emph{must-in set}.  The head $h$ is an
administrative label attached to the branch: root branches have $M=\{h\}$, and
sibling-generated branches use as head the first post-sibling vertex when such a
vertex exists.  The head records the distinguished vertex that generated the
current ordered subproblem; it is not part of the clique predicate itself.  The
\emph{coverage} of the branch is
\[
  \mathrm{Ext}(M, Q) \;=\;
  \bigl\{\, K \in \mathcal{M}(G) : M \subseteq K
  \text{ and } K \setminus M \subseteq Q \,\bigr\}.
\]
A collection of branches $\mathcal{F}$ is \emph{complete} for $S \subseteq
\mathcal{M}(G)$ if $S \subseteq \bigcup_{(M_i, h_i, Q_i) \in \mathcal{F}}
\mathrm{Ext}(M_i, Q_i)$.
\end{definition}

Intuitively, $\mathrm{Ext}(M, Q)$ is the set of maximal cliques that this
branch is responsible for finding: those that contain every vertex of $M$ and
draw their remaining members from $Q$.  The goal of the algorithm is to maintain
a complete forest $\mathcal{F}$ at every step while progressively shrinking it
by finding cliques and pruning the branches that covered them.  With this
notation in place, we can now formalize the redundancy from the introduction
and show how it leads to a principled reordering rule.

% ============================================================
\section{From Redundancy to Reordering}
\label{sec:dfs}
% ============================================================

We now turn the informal redundancy from Section~\ref{sec:introduction} into a
sequence of structural statements.  The progression is as follows.  We first
fix how a clique is reached inside an ordered branch.  We then show that once a
maximal clique has been discovered, any future clique extending one of its
proper subsets must escape it.  This escape property yields the reordering rule
for one covering clique and, in the general case, the sibling-set replacement
used by \textsc{ReorderSib}.

\subsection{Ordered DFS Within a Branch}

The core search routine for a branch is an ordered DFS.  If
$K \in \mathrm{Ext}(M,Q)$ and
\[
  K \setminus M = \{w_1,\ldots,w_t\}
  \qquad\text{with}\qquad
  \sigma(w_1) < \cdots < \sigma(w_t),
\]
then $(w_1,\ldots,w_t)$ is the \emph{canonical extension sequence} of $K$
relative to $(M,Q)$.

\begin{lemma}[Canonical reachability]\label{lem:canonical}
Let $(M,h,Q)$ be a branch and let $K \in \mathrm{Ext}(M,Q)$.  If
$K \setminus M = \{w_1,\ldots,w_t\}$ is written in increasing $\sigma$-order,
then ordered DFS can reach $K$ by choosing $w_1,\ldots,w_t$ in that order and
recursing with candidate sets
\[
  Q_j
  \;=\;
  Q \cap N^+(w_1) \cap \cdots \cap N^+(w_j),
  \qquad 1 \le j \le t.
\]
\end{lemma}

\begin{proof}
Because $K \setminus M \subseteq Q$, the first vertex $w_1$ is available in
$Q$.  Since $K$ is a clique, every $w_j$ with $j \ge 2$ is adjacent to $w_1$,
and since the sequence is increasing in $\sigma$, each such $w_j$ lies in
$N^+(w_1)$.  Hence $\{w_2,\ldots,w_t\} \subseteq Q_1$.

Applying the same argument inductively, after choosing the prefix
$w_1,\ldots,w_j$, every remaining vertex
$w_{j+1},\ldots,w_t$ is adjacent to all previously chosen vertices and has
larger $\sigma$-rank than $w_j$, so it belongs to $Q_j$.  Therefore the
ordered DFS can follow the canonical extension sequence all the way to the
state $(K,Q_t)$.  We claim that $Q_t = \emptyset$.  Indeed, if some
$u \in Q_t$ existed, then $u \in Q \subseteq N(M)$, so $u$ is adjacent to every
vertex of $M$.  Also, for each $j$, membership in $N^+(w_j)$ implies that $u$
is adjacent to $w_j$.  Thus $u$ is adjacent to every vertex of
$K = M \cup \{w_1,\ldots,w_t\}$, contradicting the maximality of $K$.
Therefore $Q_t = \emptyset$, so ordered DFS reaches $K$ as a leaf.
\end{proof}

\subsection{Unique Root Ownership}

Fix the total order $\sigma$ from Section~\ref{sec:prelim}.  The initial search
starts from one root branch per vertex.

\begin{definition}[Initial forest]\label{def:forest0}
The \emph{initial forest} is
\[
  \mathcal{F}^0
  \;=\;
  \bigl\{\, (\{v\},\, v,\, N^+(v)) : v \in V \,\bigr\}.
\]
Branch $(\{v\}, v, N^+(v))$ is the unique root branch of vertex $v$.
\end{definition}

\begin{lemma}[Unique root ownership]\label{lem:partition}
Every maximal clique $K \in \mathcal{M}(G)$ belongs to exactly one root branch
of $\mathcal{F}^0$, namely the branch seeded at its minimum-rank vertex
\[
  v_* \;=\; \arg\min_{u \in K} \sigma(u).
\]
\end{lemma}

\begin{proof}
Let $v_*$ be the minimum-rank vertex of $K$.  Every vertex
$u \in K \setminus \{v_*\}$ is adjacent to $v_*$ and satisfies
$\sigma(u) > \sigma(v_*)$, so
$K \setminus \{v_*\} \subseteq N^+(v_*)$.  Hence
$K \in \mathrm{Ext}(\{v_*\}, N^+(v_*))$.

Conversely, suppose
$K \in \mathrm{Ext}(\{u\}, N^+(u))$ for some $u \ne v_*$.  Then every other
vertex of $K$ has rank strictly greater than $\sigma(u)$, which implies
$\sigma(v_*) > \sigma(u)$, contradicting the definition of $v_*$.  Thus $K$
belongs to exactly one root branch.
\end{proof}

Lemma~\ref{lem:partition} is the source of duplicate-freedom at the root level:
every maximal clique has exactly one initial owner.

\subsection{Escape From a Known Clique}

Suppose a branch $(M,h,Q)$ is being processed and a maximal clique
$C \supseteq M$ has already been found.  The next lemma captures the structural
fact that drives reordering.

\begin{lemma}[Escape lemma]\label{lem:sibling}
Let $(M,h,Q)$ be a branch and let $\mathcal{K}$ be the set of maximal cliques
found so far.  For any
$K \in \mathrm{Ext}(M,Q) \setminus \mathcal{K}$ and any covering clique
\[
  C \in \mathcal{C}(M,\mathcal{K})
  \;:=\;
  \{\, C' \in \mathcal{K} : M \subseteq C' \,\},
\]
we have
\[
  K \cap (Q \setminus C) \ne \emptyset.
\]
\end{lemma}

\begin{proof}
Since $K \notin \mathcal{K}$, we have $K \ne C$.  Because both $K$ and $C$ are
maximal cliques, neither can properly contain the other.  Since both contain
$M$, they must differ outside $M$.  Thus $K \setminus C \ne \emptyset$.
For any vertex $v \in K \setminus C$, we also have
$v \in K \setminus M \subseteq Q$ by the definition of $\mathrm{Ext}(M,Q)$.
Hence $v \in Q \setminus C$, proving the claim.
\end{proof}

Lemma~\ref{lem:sibling} says that once $M \subseteq C$ is known, the search can
never finish inside $C$ again.  Every future maximal clique extending $M$ must
\emph{escape} $C$.

\subsection{Re-ordering After One Covering Clique}

The escape lemma immediately yields a re-ordering principle in the case of a
single covering clique.

\begin{lemma}[Single-clique re-ordering]\label{lem:reorder}
Let $(M,h,Q)$ be a branch and let $C$ be a discovered maximal clique with
$M \subseteq C$.  Write
\[
  Q_{\mathrm{out}} := Q \setminus C,
  \qquad
  Q_{\mathrm{in}} := Q \cap C.
\]
Then every clique
$K \in \mathrm{Ext}(M,Q) \setminus \mathcal{K}$ can be written as
\[
  K = M \cup S \cup T,
  \qquad
  \emptyset \ne S \subseteq Q_{\mathrm{out}},
  \qquad
  T \subseteq Q_{\mathrm{in}},
\]
where $M \cup S \cup T$ is a clique.  Consequently, any exhaustive search may
postpone all vertices of $Q_{\mathrm{in}}$ until after it has selected at least
one vertex from $Q_{\mathrm{out}}$.
\end{lemma}

\begin{proof}
Take any
$K \in \mathrm{Ext}(M,Q) \setminus \mathcal{K}$.  By
Lemma~\ref{lem:sibling}, $K$ contains at least one vertex in $Q \setminus C$.
Let
$S := K \cap (Q \setminus C)$ and $T := K \cap (Q \cap C)$.  Then
$S \ne \emptyset$, $S \subseteq Q_{\mathrm{out}}$, and
$T \subseteq Q_{\mathrm{in}}$.  Since $K$ is a clique and
$K \setminus M \subseteq Q$, we obtain
$K = M \cup S \cup T$ with the stated properties.  The search-order conclusion
follows immediately: every valid continuation must pick from $Q_{\mathrm{out}}$
before any clique completed only inside $Q_{\mathrm{in}}$ could be relevant.
\end{proof}

This is the formal version of the central design principle:
once a clique is known, its reusable vertices are pushed to the tail of the
search, while the outside vertices that enable genuinely new maximal cliques
are moved to the front.

In the motivating example, once clique $\{a,b,c\}$ is known, any later maximal
clique extending $\{a,c\}$ must first use some outside vertex such as $d$.
Re-ordering therefore means searching through $d$ before revisiting the
already-explained vertices of $\{a,b,c\}$.

\subsection{Sibling Sets Under Multiple Covering Cliques}

If several previously found maximal cliques cover the same branch, a new clique
must escape \emph{all} of them simultaneously.  This produces a natural
hitting-set problem.

\begin{definition}[Sibling set]\label{def:sibling}
For a branch $(M,h,Q)$ and discovered-clique set $\mathcal{K}$, define the
constraint family
\[
  \mathcal{H}(M,Q,\mathcal{K})
  \;=\;
  \bigl\{\, Q \setminus C : C \in \mathcal{C}(M,\mathcal{K}) \,\bigr\}.
\]
A set $S \subseteq Q$ is a \emph{sibling set} if it satisfies:
\begin{enumerate}
  \item \emph{Hitting:} $S \cap H \ne \emptyset$ for every
        $H \in \mathcal{H}(M,Q,\mathcal{K})$.
  \item \emph{Clique compatibility:} $S$ is a clique in $G$.
  \item \emph{Minimality:} no proper subset of $S$ satisfies both (1) and (2).
\end{enumerate}
We write $\mathcal{S}(M,Q,\mathcal{K})$ for the family of all sibling sets.
\end{definition}

For a sibling set $S$, the remaining common expansion vertices are
\[
  \mathrm{CE}(Q,S)
  \;=\;
  (Q \setminus S) \cap \bigcap_{v \in S} N(v).
\]

\begin{lemma}[Sibling replacement is complete]\label{lem:sibling_branch}
Let $(M,h,Q)$ be a branch and let $\mathcal{K}$ be the set of already
discovered maximal cliques.  Define
\[
  \mathcal{F}^{S}_{0}
  \;=\;
  \left\{
    (M \cup S,\, h,\, \emptyset)
    : S \in \mathcal{S}(M,Q,\mathcal{K}),\;
      \mathrm{CE}(Q,S)=\emptyset
  \right\},
\]
\[
  \mathcal{F}^{S}_{1}
  \;=\;
  \left\{
    (M \cup S \cup \{v\},\, v,\, \mathrm{CE}(Q,S) \cap N^+(v))
    : \substack{S \in \mathcal{S}(M,Q,\mathcal{K}),\\
                v \in \mathrm{CE}(Q,S)}
  \right\},
\]
and let $\mathcal{F}^{S} = \mathcal{F}^{S}_{0} \cup \mathcal{F}^{S}_{1}$.
Then $\mathcal{F}^{S}$ is complete for
$\mathrm{Ext}(M,Q) \setminus \mathcal{K}$.
\end{lemma}

\begin{proof}
Let $K \in \mathrm{Ext}(M,Q) \setminus \mathcal{K}$.  By
Lemma~\ref{lem:sibling}, $K \setminus M$ hits every set in
$\mathcal{H}(M,Q,\mathcal{K})$.  Since $K$ is a clique, $K \setminus M$ is also
clique-compatible.  Therefore $K \setminus M$ contains some inclusion-minimal
clique-compatible hitting subset $S_K$, and by construction
$S_K \in \mathcal{S}(M,Q,\mathcal{K})$.

If $K = M \cup S_K$, then $\mathrm{CE}(Q,S_K)=\emptyset$.  Otherwise some
vertex $u \in \mathrm{CE}(Q,S_K)$ would lie in $Q \setminus S_K$, would be
adjacent to every vertex of $S_K$ by definition of $\mathrm{CE}$, and would be
adjacent to every vertex of $M$ because $Q \subseteq N(M)$.  Hence $u$ would be
adjacent to every vertex of the clique $K = M \cup S_K$, contradicting the
maximality of $K$.  Therefore
$K \in \mathrm{Ext}(M \cup S_K,\emptyset)$, which is represented by the leaf
branch in $\mathcal{F}^{S}$.

Otherwise let $w$ be the minimum-rank vertex of $K \setminus (M \cup S_K)$.
Every vertex in $K \setminus (M \cup S_K \cup \{w\})$ lies in
$\mathrm{CE}(Q,S_K)$ because it belongs to $Q \setminus S_K$ and is adjacent to
every vertex of $S_K$.  By minimality of $w$, every such vertex also lies in
$N^+(w)$.  Hence
\[
  K \in \mathrm{Ext}\bigl(M \cup S_K \cup \{w\},
  \mathrm{CE}(Q,S_K) \cap N^+(w)\bigr),
\]
and the corresponding branch is present in $\mathcal{F}^{S}$.
\end{proof}

Lemma~\ref{lem:sibling_branch} is the formal sibling effect: instead of
traversing all descendants of a non-maximal subset $M$, we replace that search
space by the smaller family of branches generated from minimal escape sets.

\begin{corollary}[No sibling set, no undiscovered clique]\label{cor:no-sibling}
Let $(M,h,Q)$ be a branch.  If
$\mathcal{S}(M,Q,\mathcal{K}) = \emptyset$, then
$\mathrm{Ext}(M,Q) \setminus \mathcal{K} = \emptyset$.
\end{corollary}

\begin{proof}
If some
$K \in \mathrm{Ext}(M,Q) \setminus \mathcal{K}$ existed, then the proof of
Lemma~\ref{lem:sibling_branch} would produce an inclusion-minimal
clique-compatible hitting subset
$S_K \in \mathcal{S}(M,Q,\mathcal{K})$.  This contradicts the assumption that
$\mathcal{S}(M,Q,\mathcal{K}) = \emptyset$.
\end{proof}

% ============================================================
\section{The Core \textsc{ReorderSib} Algorithm}
\label{sec:algorithm}
% ============================================================

\subsection{Core Search Procedure}

The previous section identified the mathematical structure behind the
reordering idea.  We now package those facts into a clean core algorithm.  The
algorithm maintains a global set $\mathcal{K}$ of discovered maximal cliques
and a guard set $\mathcal{D}$ of already reported empty-candidate branches.
It starts from the initial forest $\mathcal{F}^0$ and repeatedly alternates
between two actions.

For clarity, this section presents a forest-level core algorithm: a discovered
clique reorders the residual forest in one step.  Section~\ref{sec:implementation}
then explains how the code realizes the same transformations through a more
local branch-by-branch recursion.

\textbf{Sibling phase.}  Before a non-root branch is expanded, all already found
maximal cliques covering its must-in set are collected.  The branch is then
replaced by the sibling-generated family from
Lemma~\ref{lem:sibling_branch}.  If several generated branches have the same
must-in set, they are merged by taking the union of their candidate sets; the
stored head label may then be chosen arbitrarily, since coverage depends only on
the pair $(M,Q)$.

\textbf{Enumeration phase.}  DFS is run on the current branches in order.  The
DFS branches in increasing $\sigma$-order, always recursing with
$Q \cap N^+(v)$ so that the next expansion is ordered by the new head vertex
$v$.  As soon as one maximal clique $C$ is found, the current DFS stops, $C$ is
added to $\mathcal{K}$, and the residual forest is rebuilt from the current
branch together with all later branches by enlarging each candidate set with the
vertices of $C$ that are compatible with its must-in set.  A leaf $R$ is
reported only if it is globally maximal, i.e., $N(R)=\emptyset$.  The algorithm
then recurses on the resulting residual forest, where the next sibling phase can
use $C$ as a covering clique.

This is the recursive re-ordering strategy: a discovered clique is used
immediately to move already-explained vertices to the tail and to expose the
outside vertices that can still produce new maximal cliques.

\begin{algorithm}[t]
\caption{\textsc{rCall}: sibling reduction and enumeration dispatch}
\label{alg:rcall}
\begin{algorithmic}[1]
\Require Forest $\mathcal{F} = \{(M_i,h_i,Q_i)\}$, level $\ell \ge 0$,
         global discovered set $\mathcal{K}$, guard set $\mathcal{D}$
\If{$\ell > 0$ \textbf{and} $\mathcal{F} \ne \emptyset$}
  \State $(M,h,Q) \gets (M_1,h_1,Q_1)$
  \State $\mathcal{F}_{\mathrm{rest}} \gets \mathcal{F} \setminus \{(M,h,Q)\}$
  \State $\mathcal{S} \gets \mathcal{S}(M,Q,\mathcal{K})$
  \If{$\mathcal{S} = \emptyset$}
    \State $\mathcal{F} \gets \mathcal{F}_{\mathrm{rest}}$
  \Else
    \State $\mathcal{F}_{\mathrm{new}} \gets \emptyset$
    \ForAll{$S \in \mathcal{S}$}
      \State $Q_S \gets \mathrm{CE}(Q,S)$
      \If{$Q_S = \emptyset$}
        \State $\mathcal{F}_{\mathrm{new}} \gets \mathcal{F}_{\mathrm{new}} \cup \{(M \cup S,\; h,\; \emptyset)\}$
      \Else
        \ForAll{$v \in Q_S$}
          \State $\mathcal{F}_{\mathrm{new}} \gets \mathcal{F}_{\mathrm{new}} \cup
                 \{(M \cup S \cup \{v\},\; v,\; Q_S \cap N^+(v))\}$
        \EndFor
      \EndIf
    \EndFor
    \State $\mathcal{F} \gets \mathcal{F}_{\mathrm{new}} \cup \mathcal{F}_{\mathrm{rest}}$
  \EndIf
  \State Merge branches with identical must-in sets by unioning their candidate sets and retaining an arbitrary head label
\EndIf
\ForAll{$(M_i,h_i,Q_i) \in \mathcal{F}$}
  \State $\mathrm{done} \gets \mathrm{false}$
  \State \Call{Enumerate}{$M_i,\; h_i,\; Q_i,\; \mathcal{F},\; i,\; \ell,\;
         \mathrm{done}$}
  \If{$\mathrm{done}$} \textbf{break} \EndIf
\EndFor
\end{algorithmic}
\end{algorithm}

\begin{algorithm}[t]
\caption{\textsc{Enumerate}: ordered DFS with post-discovery reordering}
\label{alg:enumerate}
\begin{algorithmic}[1]
\Procedure{Enumerate}{$R,\; h,\; Q,\; \mathcal{F},\; i,\; \ell,\; \mathrm{done}$}
\If{$Q = \emptyset$}
  \If{$N(R) = \emptyset$ \textbf{and} $R \notin \mathcal{D}$}
    \State $\mathcal{D} \gets \mathcal{D} \cup \{R\}$
    \State $\mathcal{K} \gets \mathcal{K} \cup \{R\}$;\quad
           $\mathrm{done} \gets \mathrm{true}$
    \State $\mathcal{F}' \gets \emptyset$
    \ForAll{$(M_j,h_j,Q_j) \in \mathcal{F}$ with $j \ge i$}
      \State $Q_j' \gets \bigl(Q_j \cup (R \cap N(M_j))\bigr) \setminus M_j$
      \State $\mathcal{F}' \gets \mathcal{F}' \cup \{(M_j,h_j,Q_j')\}$
    \EndFor
    \State \Call{rCall}{$\mathcal{F}',\; \ell + 1$}
  \EndIf
  \State \Return
\EndIf
\ForAll{$v \in Q$ in increasing $\sigma$-order}
  \State \Call{Enumerate}{$R \cup \{v\},\; v,\; Q \cap N^+(v),\;
         \mathcal{F},\; i,\; \ell,\; \mathrm{done}$}
  \If{$\mathrm{done}$} \Return \EndIf
\EndFor
\EndProcedure
\end{algorithmic}
\end{algorithm}

\subsection{Correctness}

We now prove soundness, completeness, and duplicate-freedom for the core
algorithm.

\begin{theorem}[Soundness]\label{thm:sound}
Every set reported by \textsc{ReorderSib} is a maximal clique of $G$.
\end{theorem}

\begin{proof}
The algorithm reports a set $R$ only at a leaf with $Q=\emptyset$ and
$N(R)=\emptyset$.  Since every recursive extension adds only vertices adjacent
to all previously chosen vertices, the reported set $R$ is a clique.  The
condition $N(R)=\emptyset$ is exactly the definition of global maximality.
Hence every reported set is a maximal clique.
\end{proof}

\begin{theorem}[Completeness]\label{thm:complete}
\textsc{ReorderSib} reports every maximal clique in $\mathcal{M}(G)$.
\end{theorem}

\begin{proof}
We maintain the invariant that at the start of every call
$\textsc{rCall}(\mathcal{F},\ell)$, the forest $\mathcal{F}$ is complete for
$\mathcal{M}(G) \setminus \mathcal{K}$.

\emph{Base case.}  Initially, $\mathcal{K}=\emptyset$ and
$\mathcal{F}=\mathcal{F}^0$.  By Lemma~\ref{lem:partition},
$\mathcal{F}^0$ is complete for $\mathcal{M}(G)$.

\emph{Sibling phase.}  When $\ell>0$, the first branch $(M,h,Q)$ is replaced by
the family from Lemma~\ref{lem:sibling_branch}.  If
$\mathcal{S}(M,Q,\mathcal{K})=\emptyset$, then Corollary~\ref{cor:no-sibling}
shows that no undiscovered maximal clique remains in $\mathrm{Ext}(M,Q)$, so
the branch may be dropped.  Otherwise the
replacement family is complete for $\mathrm{Ext}(M,Q) \setminus \mathcal{K}$,
while every other branch is left unchanged.  Hence completeness of the whole
forest is preserved.

\emph{Existence of the next report.}  Since $\mathcal{F}$ is complete for
$\mathcal{M}(G)\setminus\mathcal{K}$, if this set is non-empty then some branch
$(M,h,Q)$ covers a maximal clique $K \in \mathcal{M}(G)\setminus\mathcal{K}$.
By Lemma~\ref{lem:canonical}, ordered DFS on that branch can reach $K$ along
its canonical extension sequence and reaches it as a leaf.  Since $K$ is a
maximal clique, the reporting test $N(K)=\emptyset$ is satisfied at that leaf.
Because the DFS is finite and exhaustive within each branch, the algorithm
eventually reports some maximal clique $C \in
\mathcal{M}(G)\setminus\mathcal{K}$.

\emph{Forest update after a report.}  Suppose DFS reports a maximal clique $C$.
Let $\mathcal{K}' = \mathcal{K} \cup \{C\}$.  Consider any
$K \in \mathcal{M}(G) \setminus \mathcal{K}'$.
Since branch $i$ is the first branch in the current call that reports a clique,
every earlier branch $1,\ldots,i-1$ has already been searched exhaustively and
contains no clique in $\mathcal{M}(G)\setminus\mathcal{K}'$.  It therefore
suffices to consider branches $j \ge i$.

By the induction hypothesis, $K$ was covered by some pre-update branch
$(M_j,h_j,Q_j)$ with $j \ge i$.  In the residual forest, this branch is kept
with updated candidate set
\[
  Q_j' = \bigl(Q_j \cup (C \cap N(M_j))\bigr) \setminus M_j.
\]
Since $Q_j \subseteq Q_j'$, every clique that was covered by $(M_j,h_j,Q_j)$ is
still covered by $(M_j,h_j,Q_j')$.  In particular,
$K \in \mathrm{Ext}(M_j,Q_j')$.  Hence the residual forest remains complete for
$\mathcal{M}(G) \setminus \mathcal{K}'$.

By induction, every maximal clique is eventually reported.
\end{proof}

\begin{theorem}[Duplicate-freedom]\label{thm:unique}
\textsc{ReorderSib} reports each maximal clique at most once.
\end{theorem}

\begin{proof}
There are three possible sources of duplication.

\emph{Within one DFS branch.}  In \textsc{Enumerate}, once the DFS chooses a
head vertex $v$, all later recursion is restricted to $N^+(v)$.  Hence the
vertices added after the head are visited in strictly increasing $\sigma$-order,
so a fixed branch cannot report the same maximal clique twice.

\emph{Across sibling-generated branches with the same must-in set.}  The
algorithm explicitly merges such branches before DFS begins, so they cannot lead
to duplicate reports.

\emph{Across different later branches in the same forest.}  Suppose a maximal
clique $K$ could be reported again after it has already been discovered.  This
cannot happen because every reported clique is inserted into the guard set
$\mathcal{D}$ immediately.  Any later leaf equal to $K$ is therefore discarded
without being reported.

Hence no maximal clique is reported more than once.
\end{proof}

% ============================================================
\section{Implementation Realization}
\label{sec:implementation}
% ============================================================

The previous section described the theorem-level algorithm used for the
correctness proof.  We now return to the concrete implementation in
\texttt{helpers.cpp}.  The implementation realizes the same reordering idea,
but with explicit branch-management data structures and more aggressive
operational shortcuts.

\paragraph{Branch representation.}
The implementation stores each branch as a pair of arrays:
\texttt{mustin[i]} for the must-in set $M$ and \texttt{expandTo[i]} for the
candidate set $Q$.  The code does not store the abstract head $h$ explicitly.
For ordinary ordered branches, the controlling vertex is recovered from
\texttt{mustin[i].back()}, while sibling-generated branches are tracked through
\texttt{fullSkipCheck} and whole-branch containment tests rather than through a
separate head field.  The recursive driver \texttt{rCall} manages the current
local branch family, while \texttt{enumerate} performs the DFS and triggers the
next reordering step when it finds a clique.

\paragraph{Covering-clique retrieval.}
All discovered cliques are stored globally in \texttt{allCliques}.  To avoid
scanning that list naively, the code maintains a two-tier index
\[
  \mathcal{I}[v][\ell']
\]
whose entries are clique IDs containing vertex $v$ and discovered at level
$\ell'$.  The function \texttt{collectCoveringCliques} chooses the least
frequent vertex of the must-in set as a seed, filters on that vertex, and then
verifies containment of the whole must-in set.  This is the implementation
counterpart of the set $\mathcal{C}(M,\mathcal{K})$.

\paragraph{Sibling effect.}
Given a branch $(M,Q)$ and its covering cliques, the implementation constructs
the constraint family $\{Q \setminus C\}$ and solves the corresponding
clique-compatible hitting-set problem in
\texttt{generateSiblingSetsFromCliques}.  The function
\texttt{commonExpand} computes the common expansion
$\mathrm{CE}(Q,S)$.  The local forest created by \texttt{rCall} is exactly the
branch family described in Lemma~\ref{lem:sibling_branch}, except that the code
materializes the family one local subproblem at a time and then merges any two
generated branches with the same must-in set by unioning their candidate sets.

\paragraph{Ordered DFS and recursive reordering.}
Inside \texttt{enumerate}, recursion on a chosen vertex $v$ uses
\texttt{adjList2[v]}, which is the stored forward neighborhood $N^+(v)$ under
the fixed vertex order.  When a clique is found, the implementation stops the
current DFS immediately, inserts the clique into the global discovered set, and
rebuilds the residual search space before making the next recursive
\texttt{rCall}.  Unlike the forest-level presentation in
Section~\ref{sec:algorithm}, the code launches one recursive call per surviving
reordered branch.  This is an implementation choice, not a different search
principle: known vertices are still pushed to the tail, and new escape vertices
are still brought to the front.

\paragraph{Eager branch subsumption in the implementation.}
The clean core algorithm keeps the current branch and all later branches in the
residual forest, then lets the next sibling phase use the newly reported clique
as a covering clique.  The implementation is more aggressive: it may delete a
branch immediately when cheaper operational tests indicate that the remaining
branch space has already been explained.  Ordinary branches use the last
mandatory vertex.  Sibling-generated branches additionally use
\texttt{fullSkipCheck} together with
\texttt{branchSpaceInsideClique}.  These are implementation-level accelerations
of the same reordering idea.

\paragraph{Conceptual maximality guard.}
The core algorithm in Section~\ref{sec:algorithm} uses the explicit predicate
$N(R)=\emptyset$ when a DFS leaf is reached.  The implementation does not
perform that test as a separate subroutine.  Instead, it relies on the
reordered frontier, branch deletion, and candidate pulling to ensure that
non-maximal leaves are operationally suppressed before they would be reported.
The explicit guard is part of the clean theorem-level specification.

\paragraph{Duplicate guards and aggressive branch management.}
The code keeps an explicit set \texttt{claimedEmptyBranches} to block repeated
reports from empty-candidate branches, mirroring the guard set $\mathcal{D}$ in
Algorithm~\ref{alg:enumerate}.  It also uses \texttt{fullSkipCheck} to mark
branches produced by sibling replacement, allowing the implementation to decide
whether a whole branch space is already contained in a newly found clique.
These devices are part of the concrete implementation, not part of the abstract
mathematical interface of a branch.

\paragraph{Exact hitting-set solvers.}
Two exact solvers are implemented.

\textit{Backtracking branch-and-bound.}  A DFS grows a current set
$S_{\mathrm{cur}} \subseteq Q$ in increasing order.  As soon as
$S_{\mathrm{cur}}$ hits all constraints and remains a clique, the branch stops:
no proper superset can be minimal.  A minimal-by-inclusion filter is applied to
the resulting solutions.

\textit{Bitmask exact solver.}  When the number of constraints is at most $63$,
the coverage of a candidate vertex is stored as a bitmask.  The DFS then uses
constant-time coverage updates and three exact reductions: unit propagation,
constraint subsumption, and fail-first ordering of constraints by size.  A
fourth pruning rule skips any candidate that contributes no uncovered
constraint, since such a candidate cannot participate in a minimal hitting set
at that point.  The solver maintains a live minimal-by-inclusion solution set,
so no separate dominance pass is required after the search.

\paragraph{Pruning reductions.}
The implementation also applies additional structural reductions before the
solver is invoked.  Dominance pruning is exact: it removes any covering clique
$C_i$ that is a proper subset of another covering clique $C_j$, because the
tighter constraint $Q \setminus C_j$ already implies $Q \setminus C_i$.  The
implementation also uses a level filter to limit covering-clique retrieval to
recent recursive levels.  This filter is an engineering device of the concrete
implementation and is not part of the theorem-level core proof.

\bibliographystyle{ACM-Reference-Format}
\bibliography{sample}

\end{document}
\endinput
